\section{Limitations and Future Directions}
\label{sec:limitations}

Nine constraints bound our findings. First, single-turn evaluation leaves multi-turn loops unexamined. Second, grounding varies across benchmarks rather than within them: math answers were withheld to contrast execution against unaided reasoning. Third, our 15-input fuzzing audit cannot prove code bug-free, so overrides remain unconfirmed defects. Fourth, in math and science, regex extraction and GPQA label noise make ungrounded catch rates ($26.4\%$ in math) empirical lower bounds. Fifth, style probes tested two models with 100\% baseline false approval, risking regression to the mean without altering code style. Sixth, the judge pool spans four open-weight models from one provider, omitting proprietary systems (e.g., o1). Seventh, juries and sweeps were benchmarked on $n = 150$ HumanEval+ items ($\text{temperature}=0$), which cannot rule out contamination. Eighth, single-sentence attribution is a weak manipulation for stateless models. Ninth, instructing verifiers not to re-solve from scratch may encourage belief persistence by preventing independent re-derivation. Future work includes formal provers, retrieval, and decorrelating training regimes.
